% Generated draft: requires mathematical and source review.
\documentclass[11pt]{article}
\usepackage[margin=25mm]{geometry}
\usepackage{fontspec}
\setmainfont{Latin Modern Roman}
\usepackage{amsmath,amssymb,mathtools,unicode-math}
\setmathfont{Latin Modern Math}
\usepackage{enumitem,needspace,xurl,hyperref}
\setlength{\parindent}{0pt}
\setlength{\parskip}{5pt}
\setlength{\emergencystretch}{3em}
\newcommand{\N}{\mathbb{N}}
\newcommand{\Z}{\mathbb{Z}}
\newcommand{\Q}{\mathbb{Q}}
\newcommand{\R}{\mathbb{R}}
\newcommand{\C}{\mathbb{C}}
\newcommand{\F}{\mathbb{F}}
\newcommand{\A}{\mathbb{A}}
\newcommand{\PP}{\mathbb{P}}
\newcommand{\E}{\mathbb{E}}
\newcommand{\eps}{\varepsilon}
\newcommand{\dd}{\,\mathrm{d}}
\begin{document}
\section*{UnsolvedMath ID: 30000296}
\textbf{Source problem code:} catalogue entry
\subsection{Definitions and mathematical statement}
Let $Z=\Gamma\backslash X$ be the noncompact Hermitian locally symmetric space intended by the source, with $X$ a connected Hermitian symmetric space and $\Gamma$ a discrete group acting by holomorphic isometries. For a smooth quotient one ordinarily requires a free properly discontinuous action; any orbifold interpretation must be specified if torsion is allowed. Let $\mathcal M(Z)$ be its field of global meromorphic functions, over the constant field $\C$.

Determine
\[
 \operatorname{trdeg}_{\C}\mathcal M(\Gamma\backslash X).
\]
Here the transcendence degree is the cardinality of a transcendence basis, or equivalently the maximal size of an algebraically independent family of global meromorphic functions over $\C$. A finite family $f_1,\ldots,f_r$ is algebraically independent when
\[
 P(f_1,\ldots,f_r)\not\equiv0
 \quad\text{for every nonzero }P\in\C[t_1,\ldots,t_r].
\]
The request is to determine this invariant for the intended class of pairs $(X,\Gamma)$; the extraction supplies no proposed value or criterion.

The source does not specify finite covolume, arithmeticity, rank, irreducibility, torsion conditions, or compactification data. It also does not impose automorphic growth conditions on meromorphic functions. Such restrictions cannot be added silently. In particular, noncompactness alone does not justify importing a compact-complex-space bound on transcendence degree, so neither finiteness nor equality with $\dim_{\C}Z$ is asserted here.

The literature assessment additionally discusses a Helgason Poisson-transform conjecture and says it was proved by Kashiwara and collaborators. That conjecture is not stated in the supplied mathematical question. Its alleged resolution must not be treated as a solution of the separate meromorphic-function-field problem, which the assessment calls the broader Siegel transcendence-degree problem.
\subsection{Short English statement}
The question asks how many algebraically independent global meromorphic functions exist on the specified noncompact Hermitian locally symmetric quotient. The quotient class and a proposed formula are missing; a separate solved Poisson-transform question mentioned in the metadata does not answer this target.
\subsection{Sources}
\begin{itemize}
\item Local UnsolvedMath registry, permanent ID 30000296, code catalogue entry, ``Meromorphic Function Fields on Locally Symmetric Spaces.'' Proposed year 2005; proposer unspecified. The assessment calls the meromorphic-function question a Siegel transcendence-degree problem.
\item \emph{Analysis and Geometric Singularities} (2006), Oberwolfach Reports, pp.\ 2187--2240, DOI 10.4171/owr/2005/39; original context \url{https://doi.org/10.4171/owr/2005/39}. The displayed report year 2006 differs from the DOI and proposed year 2005.
\item Supplied reference labeled ``modern boundary-value treatment'': \url{https://arxiv.org/abs/1807.07131}. Its relevance to the meromorphic-field question, rather than the Poisson-transform topic, is not specified.
\item Statement status corrected\_verified; the supplied verification says quotient-space and function-field notation was restored without changing the task.
\item Literature triage checked 2026-08-21: partially solved. It attributes a proof of the Helgason conjecture to Kashiwara and collaborators but says the broader Siegel problem is not solved in full generality. This bundled assessment requires source and status verification. Created 2026-08-25; updated 2026-09-16.
\item Registry curation by Ulam AI / UnsolvedMath Contributors, CC BY 4.0; underlying report and literature provenance remain as supplied. Unreviewed statement-only draft.
\end{itemize}
\end{document}
